%!TEX root = ../../dokumentation.tex

\chapter{Einleitung}
Erste Erwähnung eines Akronyms wird als Fußnote angezeigt. Jede weitere wird
nur verlinkt: \acf{AGPL}. % \cite{fsf:2007}

Verweise auf das Glossar: \gls{Glossareintrag}, \glspl{Glossareintrag}

Nur erwähnte Literaturverweise werden auch im Literaturverzeichnis gedruckt:
% \cite{baumgartner:2002}, \cite{dreyfus:1980}

Meine erste Fußnote\footnote{Ich bin eine Fußnote}

\begin{wrapfigure}{r}{.4\textwidth}
\centering
\includegraphics[height=.35\textwidth]{logo.png}
\vspace{-15pt}
\caption{Das Logo der Musterfirma\footnotemark}
\end{wrapfigure}
%Quelle muss in Fußnote stehen (da sonst aufgrund eines Fehlers nicht kompiliert
% wird)
% \footnotetext{aus \cite{mustermann:2012}}
Looking for the one superhero comic you just have to read. Following the antics and adventures of May Mayday Parker, this Spider-book has everything you could want in a comic--action, laughs, mystery and someone in a Spidey suit. Collects Alias \#1-28, What If. Jessica Jones had Joined the Avengers. In her inaugural arc, Jessicas life immediately becomes expendable when she uncovers the potentially explosive secret of one heros true identity. In her inaugural arc, Jessicas life immediately becomes expendable when she uncovers the potentially explosive secret of one heros true identity.

Once upon a time, Jessica Jones was a costumed super-hero, just not a very good one. First, a story where Wolverine and Hulk come together, and then Captain America and Cable meet up. In a city of Marvels, Jessica Jones never found her niche. The classic adventures of Spider-Man from the early days up until the 90s. Looking for the one superhero comic you just have to read. In her inaugural arc, Jessicas life immediately becomes expendable when she uncovers the potentially explosive secret of one heros true identity.

Motivation:
Problemstellung:
Zielsetzung: Aus der zuvor erarbeiteten Problemstellung ergibt sich als Ziel dieser Arbeit…

\section{Geschäftlicher Kontext}

In dieser Arbeit wird die Migration einer lokal laufenden Anwendung aus dem Projektmanagement zur Rechnungsstellung in die Cloud untersucht.
Zweck der Anwendung ist das Einsammeln sogenannter Timesheets von Projektmitarbeitern aus einem Cloud Speicher. Diese Timesheets dokumentieren,
wie viel Arbeitszeit ein Mitarbeiter für die verschiedenen Aufgaben in einem Projekt verbracht hat um daraus präzise Rechnungen erstellen zu können.

Das Tool wird bereits heute von mehreren Personen im Projekt verwendet und soll in Zukunft auch für andere Projekte eingesetzt werden. Aus diesem
Grund wird die Migration dieser Anwendung in die Cloud untersucht.

Die durchgeführte Migration soll die Anwendung daher global verfügbar machen und vorallem den Aufwand aber auch gegebenenfalls Kosten reduzieren.
Diese Anwendung muss bisher vom Anwender aus GitHub heruntergeladen und lokal ausgeführt werden.
Zur lokalen Ausführung sind einige Anpassungen in Konfigurationsdateien notwendig, da sich das ausführende System ändern kann.
Darüber hinaus muss die Anwendung nach jeder Aktualisierung manuell erneut heruntergeladen werden.
Das zur Verfügung stellen in der Cloud soll die Benutzung der Anwendung erleichtern und es auf lange Sicht ermöglichen, dass diese auch außerhalb der Projektumgebung verwendet werden kann.
\section{Einordnung dieser Arbeit in den wissenschaftlichen Kontext}

Grundlage für diese Arbeit bilden die zu dem Thema Cloud Migration veröffentlichten Publikationen, welche die Herausforderungen der Cloud Migration in Anbetracht der Wahl einer passenden Architektur erarbeiten. Dazu zählen unter anderem auch die Publikationen von Pahl et al., Carrasco et al. (2018) und Kazanavičius et al. (2019), die auf die Migration von \textit{legacy} Anwendungen in die Cloud eingehen und damit verbundene Herausforderungen aufzeigen. Dort wird unter anderem die Auswahl einer passenden Architektur \cite[Vgl.][S. 14]{Pahl}, der für das Unternehmen zutreffenden Migrationsstrategie \cite[Vgl.][S. 4]{Kazanavicius2019} und die Analyse potenzieller Vor- und Nachteile, von zum Beispiel der Microservice-Architektur \cite[Vgl.][S. 3]{Carrasco2018}, als Herausforderung erarbeitet.
\section{Kritische Auswahl der Forschungsmethoden}
\label{sec:auswahl_forschungsmethoden}

% Zielsetzung
% Literaturrecherche
% -hierzu zunächst aktuellen Stand der Forschung untersuchen
%  ->Thema x,y,z
% -Methodik beschreiben
% Anforderungsanalyse
% -non-functional Requirements
% -was muss beachtet werden für Cloud Anwendung
% -Anforderungen Cloud
% Use-Case Analyse
% -Werden Use-Cases immer noch erfüllt?
%  ->zufriedenstellend?
% -Architekturelle Änderungen (Client -> Web)
% Prototyping
% Evaluation

Um das Ziel dieser Arbeit zu erreichen wird zunächst der aktuelle Stand der wissenschaftlichen Forschung untersucht. Besonderer Fokus wird hierbei auf die Untersuchung der Herausforderungen der Cloud Migration, bevor die Durchführung eines Migrationsansatzes betrachtet wird.

Um die aktuellen Herausforderungen herauszuarbeiten wird eine Literaturrecherche nach Döring/Bortz 2016 durchgeführt. Dazu werden einige Suchbegriffe zur systematischen Durchsuchung wissenschaftlicher Datenbanken festgelegt \cite[Vgl.][S. 158]{Doering2016}. Als Suchmaschine wurde in diesem Fall hauptsächlich Google Scholar genutzt, da hierüber auch auf weitere Datenbanken verwiesen wird. Darüber hinaus wurde vor allem noch die Datenbank des Verlages Springer verwendet, aber vereinzelt auch noch einige weitere wie arXiv oder ACM.

Zur Recherche wurden die folgenden Begriffe verwendet: \glqq{Cloud Migration}\grqq{}, \glqq{Cloud Computing}\grqq{}, \glqq{Cloud}\grqq{}, \glqq{Herausforderungen Cloud Migration}\grqq{}

Diese werden als \textbf{\glqq{primäre Suchbegriffe}\grqq{}} \cite[S. 158]{Doering2016} festgelegt. Die Suche wurde darüber hinaus nach dem Schneeballsystem durchgeführt, was bedeutet, dass die Quellen der verwendeten Veröffentlichungen zur weiteren Recherche verwendet werden \cite[Vgl.][S. 160]{Doering2016}. Aus den Quellen werden außerdem weitere Begriffe als \textbf{\glqq{sekundäre Suchbegriffe}\grqq{}} \cite[S. 158]{Doering2016} festgelegt und zur erweiterten Suche verwendet. Zu diesen gehören unter anderem \glqq{Migration zu PaaS}\grqq, \glqq{Cloud Native}\grqq{} und \glqq{Refactoring}\grqq{}.

Außerdem wird das \textit{Prototyping} wie in Wilde/Hess 2007 und Heinrich 2011 beschrieben, zur Umsetzung des Anwendungsbeispiels und eine Use-Case Modellierung und Analyse dieser in Anbetracht der entwickelten Anwendung durchgeführt.